\section{The Controlled Replication Pipeline}
\label{sec:pipeline}

Figure~\ref{fig:architecture} summarizes the eight-step pipeline that
produces every quantity in \S\ref{sec:results}. Only two steps involve the
language model at all -- extraction (Step 1) and signal-code generation
(Step 3) -- and both are bounded, in the sense made precise below, so that
no empirical parameter ever originates from the model itself.

\begin{figure}[htbp]
  \centering
  \fbox{\parbox{0.9\linewidth}{\centering \vspace{4cm} F1: pipeline architecture diagram (TikZ or exported PDF) \vspace{4cm}}}
  \caption{The eight-step controlled replication pipeline. The LLM appears
    only in Step 1 (extraction) and Step 3 (\texttt{compute\_signal}
    generation).}
  \label{fig:architecture}
\end{figure}

\subsection{MethodSpec as the sole carrier of method}

Step 1 extracts a structured \emph{MethodSpec} from the paper's text alone
-- signal definition, universe, timing, and portfolio-construction fields,
each tagged with the source sentence it was read from. No downstream step
ever reads the paper directly; every later step's notion of "what the paper
says" is mediated entirely through this one artifact, which makes the
extraction auditable independently of everything that follows. Step 2 is a
review gate: a human confirms or corrects any field the extractor flags as
low-confidence or internally inconsistent, and the MethodSpec is marked
\texttt{codegen\_ready} only once every required field has either direct
textual support or an explicit human correction. This gate is also where
the RQ2 measurements originate (\S\ref{sec:intro}): every human correction
is recorded as an evidence entry on the field it touches, so the fraction of
high-impact fields the agent resolved unassisted is a byproduct of running
Step 2, not a separately constructed measurement.

\subsection{Fixed configuration menu and clamping}

Once a MethodSpec is codegen-ready, \texttt{registry.build\_config} turns it
into an executable configuration by resolving every relevant field against
a fixed menu of allowed values (weighting rule, breakpoint convention,
rebalance frequency, holding period, universe filters, estimator; the full
menu is reproduced in Appendix~\ref{app:config-menu}). A value the MethodSpec
specifies that falls outside this menu is not passed through -- it is
deterministically clamped to a documented default, and the clamp is logged
in \texttt{config["defaults\_applied"]} rather than silently dropped. This
is the mechanism that makes $\ImplSpace{P}$ (\S\ref{sec:framework})
something one can actually enumerate: because the menu is finite and
explicit, every configuration this paper ever runs is a labeled point in a
known space, not an unbounded piece of hand-written logic.

\subsection{Sandbox validation and repair-then-refreeze}

Step 4 validates the generated \texttt{compute\_signal} plugin in a sandbox
-- syntax and schema checks, plus a static scan for look-ahead leakage --
before it is allowed to run against real data. If a technical repair to the
plugin is needed mid-batch, the entire batch of tracks for that factor is
re-executed under the repaired code with repair disabled
(\texttt{max\_refreeze\_attempts} bounds how many times this can happen),
so that no two tracks reported for the same factor are ever silently
compared under different code. This guarantees that the frozen signal $\SA$
underlying $\Aref$, $\Acz$, and $\Ahxz$ is, by construction, the exact same
artifact across all three tracks -- the identity in
Equation~\eqref{eq:three-term} would not otherwise hold exactly.

\subsection{Evidence store and hashing}

Every executed track is recorded as a \texttt{RunRecord}: its configuration
hash, its plugin's code hash, and its resulting metrics, stored alongside
both an artifact hash (of the physical output file) and a semantic hash (of
its canonicalized content, invariant to representation details such as
column order or float formatting). Step 7 assembles these records for a
factor into a single evidence bundle -- the field-level configuration diff,
the gap decomposition, and the identification level of each comparison
(\S\ref{sec:design}) -- which is the only input the next step is permitted
to read.

\subsection{The Step 8 diagnosis layer}

Step 8 is an optional, LLM-assisted explanation layer over an already-
computed evidence bundle. It is deliberately not permitted to alter, or even
see, any number the deterministic engine produced: each claim it drafts must
cite one or more \texttt{evidence\_keys} already present in the bundle from
Step 7, a validator rejects any claim that cites a key not present in the
bundle or that contains a digit of its own, and every output is tagged
\texttt{status: "llm\_assisted\_proposal"} rather than presented as a
finding. This confines the language model, even in the one place it is
permitted to comment on results, to drafting attribution language for
numbers it did not compute and cannot alter.

